\chapterbackmatter{Solutions to Supplementary Exercises}

\begin{enumerate}
  \SupplementarySolution{1}{Curvature from Parallel Transport}
  \begin{enumerate}
    \item[(a)]
    Acting on a test field,
    \[
      [D_\mu,D_\nu]\psi
        =-i\left(\partial_\mu A_\nu-\partial_\nu A_\mu
          -i[A_\mu,A_\nu]\right)\psi .
    \]
    Thus
    \[
      F_{\mu\nu}
        =\partial_\mu A_\nu-\partial_\nu A_\mu-i[A_\mu,A_\nu],
    \]
    which in components is Eq.~\eqref{eq:15.1.13}.  A finite gauge
    transformation is defined by requiring
    \(D'_\mu\psi'=UD_\mu\psi\) with \(\psi'=U\psi\); hence
    \(D'_\mu=UD_\mu U^{-1}\).  It follows without further expansion that
    \[
      [D'_\mu,D'_\nu]
        =U[D_\mu,D_\nu]U^{-1}
        =-iUF_{\mu\nu}U^{-1}.
    \]
    Comparing this with
    \([D'_\mu,D'_\nu]=-iF'_{\mu\nu}\) gives
    \(F'_{\mu\nu}=UF_{\mu\nu}U^{-1}\).
    \item[(b)]
    The transporter on a short displacement \(u^\mu\) from \(z\) is
    \[
      U(z+u,z)
        =1+i u^\mu A_\mu(z)
         +\frac{i}{2}u^\mu u^\nu\partial_\mu A_\nu(z)
         -\frac12u^\mu u^\nu A_\mu(z)A_\nu(z)+O(u^3).
    \]
    Multiply the corresponding expressions for the four oriented sides in
    their traversal order.  The terms linear in \(a\) or \(b\) cancel.  The
    derivative terms give
    \(i(\partial_\mu A_\nu-\partial_\nu A_\mu)a^\mu b^\nu\), while the
    products whose order differs on the two routes give
    \([A_\mu,A_\nu]a^\mu b^\nu\).  Therefore
    \[
      U_\square
        =1+i\bigl(\partial_\mu A_\nu-\partial_\nu A_\mu
          -i[A_\mu,A_\nu]\bigr)a^\mu b^\nu+O(\ell^3)
        =1+iF_{\mu\nu}a^\mu b^\nu+O(\ell^3).
    \]
    Replacing the ordered product by an ordinary exponential would
    symmetrize the matrix factors and lose the commutator; path ordering is
    precisely what retains the non-Abelian part of the curvature.
  \end{enumerate}

  \SupplementarySolution{2}{Yang--Mills Equation and Covariant Conservation}
  \begin{enumerate}
    \item[(a)]
    Applying a covariant divergence to the field equation gives
    \[
      D_\nu J^\nu=D_\nu D_\mu F^{\mu\nu}
        =\frac12[D_\nu,D_\mu]F^{\mu\nu}.
    \]
    In the adjoint representation,
    \([D_\nu,D_\mu]X=-i[F_{\nu\mu},X]\), up to the convention for where
    \(g\) is absorbed.  Therefore
    \[
      D_\nu J^\nu
        =-\frac{i}{2}[F_{\nu\mu},F^{\mu\nu}]=0.
    \]
    The last contraction vanishes because interchanging the two field
    strengths reverses the Lie bracket while leaving the summed Lorentz
    indices unchanged.  Thus the Yang--Mills equation requires the
    covariant, rather than ordinary, conservation law
    \[
      \boxed{D_\nu J^\nu=0.}
    \]
    \item[(b)]
    Since \(\Psi'=U\Psi\) and
    \(\overline\Psi'=\overline\Psi U^{-1}\),
    \[
      J'^\nu
        =g\overline\Psi U^{-1}\gamma^\nu T^aU\Psi\,T^a
        =UJ^\nu U^{-1},
    \]
    where the final equality is the matrix form of the adjoint action on
    the components.

    The transformed covariant derivative satisfies
    \(D'_\mu=UD_\mu U^{-1}\), while
    \(F'^{\mu\nu}=UF^{\mu\nu}U^{-1}\).  Acting on the curvature therefore
    gives
    \[
      D'_\mu F'^{\mu\nu}
        =U(D_\mu F^{\mu\nu})U^{-1}.
    \]
    Both sides of the field equation transform homogeneously in the
    adjoint representation, so its covariance follows from the independent
    transformations of the matter current and curvature.
  \end{enumerate}

  \SupplementarySolution{3}{Equal Curvature Is Not Gauge Equivalence}
  \begin{enumerate}
    \item[(a)]
    With \(F_{12}=\partial_1A_2-\partial_2A_1+[A_1,A_2]\), the first
    connection gives
    \[
      F_{12}=a^2[T_1,T_2]=a^2T_3.
    \]
    The second gives
    \(F'_{12}=\partial_1(a^2x^1T_3)=a^2T_3\).
    Nevertheless their covariant derivatives of the curvature differ:
    \[
      D_1F_{12}=a^3[T_1,T_3]\ne0,\qquad
      D_2F_{12}=a^3[T_2,T_3]\ne0,
    \]
    while \(D'_1F'_{12}=D'_2F'_{12}=0\).  Hence, for example,
    \[
      \operatorname{tr}\bigl(D_\rho F_{12}D^\rho F_{12}\bigr)
    \]
    distinguishes the connections.  This scalar is gauge invariant, so the
    two connections cannot be related by a gauge transformation.  In an
    Abelian theory equal curvatures imply \(d(A-A')=0\).  The Poincaré lemma
    makes \(A-A'=d\lambda\) on a simply connected contractible patch, which
    is exactly an Abelian gauge transformation.
  \end{enumerate}

  \SupplementarySolution{4}{Gauge Transport and Curvature}
  \begin{enumerate}
    \item[(a)] Endpoint factors cancel around the quadrangle, so its
      ordered product transforms by conjugation at the base point.  Applying
      the short-segment expansion on each edge, including the translated
      connection, gives
      \[
        \mathcal W(x)=1+i\epsilon^2m^\mu n^\nu
        \bigl(\partial_\mu A_\nu-\partial_\nu A_\mu
              i[A_\mu,A_\nu]\bigr)+O(\epsilon^3).
      \]
      The bracket is \(F_{\mu\nu}\), up to the orientation and covariant
      derivative conventions.
    \item[(b)] Commutativity removes path ordering in the Abelian case.
      Thus, for \(C=\partial\Sigma\),
      \[
        U[C]=\exp\left(i\oint_C A\right)
          =\exp\left(i\int_\Sigma F\right).
      \]
      Two surfaces give the same result when their union bounds a region
      with \(dF=0\); magnetic sources and nontrivial topology are the
      standard qualifications.
    \item[(c)] Put
      \(I(t,s)=\int_s^t d\tau\,\dot y^\mu A_\mu(y(\tau))\).
      Differentiating \(U=e^{iI}\) gives
      \[
        \partial_tU=i\dot y(t)\cdot A(y(t))U,\qquad
        \partial_sU=-iU\dot y(s)\cdot A(y(s)).
      \]
    \item[(d)] The ordered expansion is
      \[
        U(t,s)=1+\sum_{r\geq1}i^r
        \int_{s<\tau_r<\cdots<\tau_1<t}
        d^r\tau\,A(\tau_1)\cdots A(\tau_r).
      \]
      Differentiating the upper boundary inserts \(A(t)\) on the left;
      differentiating the lower boundary inserts \(A(s)\) on the right
      with a minus sign.  These are precisely the two endpoint equations.
    \item[(e)] Let
      \(U'(t,s)=V(y(t))U(t,s)V^{-1}(y(s))\).  The upper endpoint equation
      is preserved provided
      \[
        A'_\mu=VA_\mu V^{-1}-i(\partial_\mu V)V^{-1}.
      \]
      The lower equation yields the same law and the correct right endpoint
      multiplication.  Hence the complete pair of transport equations is
      gauge covariant.
  \end{enumerate}

  \SupplementarySolution{5}{Nonlinear Lorenz and Axial Gauges}
  \begin{enumerate}
    \item[(a)] The infinitesimal Abelian gauge change gives
      \[
        \delta G=(\partial^2+\zeta A^\mu\partial_\mu)\lambda,
      \]
      so \(M=\partial^2+\zeta A\cdot\partial\).  Exponentiating
      \(\det M\) and averaging \(\Omega\) gives
      \[
        \mathcal L_{\rm gf+gh}
        =-\frac1{2\xi}
          \left(\partial\cdot A+\frac{\zeta}{2}A^2\right)^2
          \omega^*(\partial^2+\zeta A\cdot\partial)\omega .
      \]
      Its quadratic inverse is
      \[
        D_{\mu\nu}(p)=\frac{-i}{p^2+i0}
          \left[\eta_{\mu\nu}-(1-\xi)\frac{p_\mu p_\nu}{p^2}\right],
        \qquad D_{\rm gh}(p)=\frac{i}{p^2+i0}.
      \]
      The nonlinear gauge also produces
      \(\zeta\omega^*A^\mu\partial_\mu\omega\),
      \(-\zeta(\partial\cdot A)A^2/(2\xi)\), and
      \(-\zeta^2(A^2)^2/(8\xi)\).  The field-dependent ghost operator
      prevents Abelian ghost decoupling.  All three extra vertices vanish
      when \(\zeta\to0\), restoring the usual covariant gauge.
    \item[(b)] In Yang--Mills theory
      \[
        \delta(n\cdot A)^a=(n\cdot D)^{ab}\lambda^b,
        \qquad M^{ab}=(n\cdot D)^{ab}.
      \]
      Inverting the quadratic kernel gives, in one common metric convention,
      \[
        D_{\mu\nu}^{ab}(p)=\frac{-i\delta^{ab}}{p^2+i0}
        \left[\eta_{\mu\nu}
          \frac{(\xi p^2+n^2)p_\mu p_\nu}{(p\cdot n)^2}
          \frac{p_\mu n_\nu+n_\mu p_\nu}{p\cdot n}\right].
      \]
      Reversing the sign used for the gauge kinetic operator reverses the
      last two displayed signs; multiplying by that operator is the
      convention-independent check.  The ghost inverse is \(n\cdot p\),
      and its only gauge interaction is proportional to
      \(C^{abc}\omega_a^*(n\cdot A_b)\omega_c\).  On the sharp slice
      \(n\cdot A=0\), \(M=n\cdot\partial\) is field independent and the
      ghosts may be omitted.  One must still fix residual
      \(n\cdot\partial\lambda=0\) transformations and prescribe the axial
      poles \(1/(p\cdot n)\).
  \end{enumerate}

  \SupplementarySolution{6}{BRST Nilpotency and Gauge-Fermion Fixing}
  \begin{enumerate}
    \item[(a)]
    Nilpotency on \(\omega^*\) and \(h\) is immediate.  On the ghost,
    the graded Leibniz rule gives
    \[
      s^2\omega_a
        =\frac14C_{abc}C_{bde}\omega_d\omega_e\omega_c
         -\frac14C_{abc}C_{cde}\omega_b\omega_d\omega_e.
    \]
    Anticommuting the ghosts puts this in the form
    \[
      s^2\omega_a
        \propto C_{ae[b}C_{ecd]}\omega_b\omega_c\omega_d=0,
    \]
    which is the Jacobi identity.  For the gauge field,
    \[
      s^2A_\mu=D_\mu(s\omega)+[sA_\mu,\omega].
    \]
    Substituting \(s\omega=-\tfrac12[\omega,\omega]\) and
    \(sA_\mu=D_\mu\omega\), the derivative terms cancel because the ghosts
    anticommute, while the remaining triple-ghost coefficient is again the
    Jacobi combination.  Hence \(s^2A_\mu=0\).  No equation of motion was
    used; the two ingredients are Grassmann antisymmetry, which selects
    antisymmetrized structure-constant products, and the Lie-algebra Jacobi
    identity, which kills them.
    \item[(b)]
    Put
    \(\Delta_a\equiv sf_a\).  Since \(\omega_a^*\) is odd, the graded
    Leibniz rule and \(s\omega_a^*=-h_a\) give
    \[
      s\Psi=\int d^4x\left[
        \omega_a^*\Delta_a+h_af_a+\frac{\xi}{2}h_ah_a\right].
    \]
    Thus the first term is precisely the Faddeev--Popov ghost action and the
    other two are the auxiliary form of gauge fixing.  Nilpotency immediately
    yields \(s(s\Psi)=0\), so this entire sector is BRST invariant.  The
    \(h_a\) equation,
    \(h_a=-f_a/\xi\), reduces it to
    \[
      \int d^4x\left(\omega_a^*\Delta_a-\frac{f_af_a}{2\xi}\right).
    \]
    The minus sign in the definition of \(\Psi\) is the one compatible with
    \(s\omega^*=-h\) and the left-acting graded convention; changing both
    conventions leaves the physics unchanged.
  \end{enumerate}

  \SupplementarySolution{7}{The Adjoint of \(SU(2)\)}
  The adjoint generators are
  \[
    t^A_1=\begin{pmatrix}0&0&0\\0&0&-i\\0&i&0\end{pmatrix},\quad
    t^A_2=\begin{pmatrix}0&0&i\\0&0&0\\-i&0&0\end{pmatrix},\quad
    t^A_3=\begin{pmatrix}0&-i&0\\i&0&0\\0&0&0\end{pmatrix}.
  \]
  Direct multiplication gives
  \([t^A_a,t^A_b]=i\epsilon_{abc}t^A_c\).  It also gives
  \[
    \sum_{a=1}^3t^A_at^A_a=2\,\mathbf1_3,
  \]
  so \(C_A=2\), as appropriate to spin one.  In the fundamental
  representation,
  \[
    \sum_at_at_a=\frac14\sum_a\sigma_a^2
      =\frac34\,\mathbf1_2,
  \]
  giving \(C_F=3/4\).  The different eigenvalues reflect the different
  representations, even though both sets realize the same Lie algebra.

  \SupplementarySolution{8}{The Complete Free Abelian BRST Multiplet}
  \begin{enumerate}
    \item[(a)] The Maxwell term is gauge invariant.  The remaining variation
      is
      \[
        s(h\,\partial\cdot A)+s(\omega^*\partial^2\omega)
        =h\partial^2\omega-h\partial^2\omega=0.
      \]
      Also \(s^2A_\mu=\partial_\mu s\omega=0\),
      \(s^2\omega^*=-sh=0\), and nilpotence on \(\omega,h\) is immediate.
    \item[(b)] After integrating derivatives by parts,
      \[
        S=-\frac12\int d^Dx\,\Phi^T
        \begin{pmatrix}
          -\partial^2\eta^{\mu\nu}+\partial^\mu\partial^\nu
            &\partial^\mu\\
          -\partial^\nu&-\xi
        \end{pmatrix}\Phi
        +\int d^Dx\,\omega^*\partial^2\omega .
      \]
      The opposite off-diagonal signs arise from the integration by parts
      and make the block operator transpose correctly.
    \item[(c)] Acting with the quadratic operators on normalized
      time-ordered correlators gives delta functions.  Fourier inversion
      yields, with the source's overall convention,
      \[
        \int d^Dx\,e^{-ipx}\langle\Phi(x)\Phi^T(0)\rangle
        =-\frac1{p^2}
        \begin{pmatrix}
          i\!\left[\eta_{\mu\nu}
             -(1-\xi)\dfrac{p_\mu p_\nu}{p^2}\right]&ip_\mu\\
          -ip_\nu&0
        \end{pmatrix},
      \]
      and
      \[
        \int d^Dx\,e^{-ipx}
          \langle\omega(x)\omega^*(0)\rangle=-\frac{i}{p^2}.
      \]
      Direct multiplication by the matrix in part (b) gives the identity
      block and fixes every mixed sign.
    \item[(d)] Since \(sh=0\),
      \[
        0=\langle s[h(x)\omega^*(0)]\rangle
          =-\langle h(x)h(0)\rangle .
      \]
      Likewise,
      \[
        0=\langle s[A_\mu(x)\omega^*(0)]\rangle
          =\langle\partial_\mu\omega(x)\omega^*(0)\rangle
           -\langle A_\mu(x)h(0)\rangle .
      \]
      The \(hh\) entry above is zero, while the mixed entry and the
      derivative of \(-i/p^2\) agree exactly, verifying both BRST Ward
      identities.
  \end{enumerate}

  \SupplementarySolution{9}{Lie-Algebra BRST Cohomology}
  \begin{enumerate}
    \item[(a)]
    On a generator,
    \[
      Q^2\omega^a
        =\frac16C^a{}_{e[b}C^e{}_{cd]}
          \omega^b\omega^c\omega^d .
    \]
    The brackets denote total antisymmetrization, forced by the Grassmann
    product.  Thus the Jacobi identity makes \(Q^2\omega^a=0\).  Because
    \(Q\) is an odd graded derivation, vanishing on all generators implies
    vanishing on every polynomial.  Conversely, if \(Q^2=0\), the
    coefficient of each independent cubic ghost monomial must vanish, which
    is precisely Jacobi.

    The polynomial algebra is the exterior algebra
    \(\Lambda^\bullet\mathfrak g^*\); \(Q\) is the
    Chevalley--Eilenberg differential for trivial coefficients.  Its ghost
    number is the cochain degree, and
    \[
      \ker Q/\operatorname{im}Q=H^\bullet(\mathfrak g)
    \]
    is the Lie-algebra cohomology.
    \item[(b)]
    Applying \(Q\) produces four ghosts:
    \[
      Q\Omega
        =-\frac1{12}C_{abc}C^a{}_{de}
          \omega^d\omega^e\omega^b\omega^c+\text{cyclic terms}.
    \]
    After total antisymmetrization, its coefficient is the invariant-metric
    form of the Jacobi identity,
    \[
      C_{a[bc}C^a{}_{de]}=0,
    \]
    so \(Q\Omega=0\).

    This is the Cartan three-cocycle and represents the nonzero generator of
    \(H^3(\mathfrak g)\) for each compact simple factor.  For
    \(\mathfrak{su}(2)\), choose \(C_{abc}=\epsilon_{abc}\).  Then
    \(\Omega=\omega^1\omega^2\omega^3\), apart from normalization, and
    \[
      Q\omega^1=-\omega^2\omega^3,\quad
      Q\omega^2=-\omega^3\omega^1,\quad
      Q\omega^3=-\omega^1\omega^2.
    \]
    Every quadratic monomial has zero \(Q\)-image because a ghost is repeated.
    Therefore no degree-two polynomial can map to \(\Omega\), proving
    non-exactness explicitly.  Restriction to the \(SU(2)\) subgroup
    associated with a root gives the same obstruction in the general compact
    simple case.
  \end{enumerate}

  \SupplementarySolution{10}{Auxiliary-Field Gauge Fixing}
  \begin{enumerate}
    \item[(a)] Completing the square gives
      \[
        h\,\partial\cdot A+\frac{\xi}{2}h^2
        =\frac{\xi}{2}\left(h+\frac{\partial\cdot A}{\xi}\right)^2
          -\frac{(\partial\cdot A)^2}{2\xi}.
      \]
      Translation invariance of \(Dh\) leaves a field-independent
      determinant and produces the standard covariant-gauge term.  The
      free Abelian ghost determinant is likewise field independent.
    \item[(b)] The Maxwell term is invariant because
      \(sF_{\mu\nu}=0\).  The matter transformation is an infinitesimal
      gauge transformation, so
      \(s[\bar\psi(i\slashed D-m)\psi]=0\).  Finally,
      \[
        s(h\partial\cdot A)=h\partial^2\omega,\qquad
        s[\omega^*(-\partial^2)\omega]
          =-h\partial^2\omega ,
      \]
      while \(s(h^2)=0\).  The two remaining variations cancel.
    \item[(c)] Clearly \(s^2A_\mu=s^2\omega=s^2\omega^*=s^2h=0\).
      For the fermion,
      \[
        s^2\psi=-ie[(s\omega)\psi-\omega s\psi]
          =e^2\omega^2\psi=0,
      \]
      because \(\omega\) is Grassmann odd; the same calculation applies
      to \(\bar\psi\).  Thus nilpotence is algebraic and off shell.
  \end{enumerate}

  \SupplementarySolution{11}{Ghost Number and Its Current}
  It is convenient to write the ghost density as
  \[
    \mathcal L_{\rm GH}
      =-(\partial_\mu\omega_a^*)(D^\mu\omega)_a
  \]
  with the adjoint derivative chosen to reproduce
  Eq.~\eqref{eq:15.6.13}.  Promote the phase to
  \(\theta(x)\).  The terms proportional to \(\theta\) cancel, while
  \[
    \delta\mathcal L_{\rm GH}
      =(\partial_\mu\theta)j^\mu_{\rm gh},\qquad
    j^\mu_{\rm gh}
      =i\left[\omega_a^*(D^\mu\omega)_a
       -(\partial^\mu\omega_a^*)\omega_a\right],
  \]
  up to an overall sign convention for the Noether charge.  Taking the
  divergence and using
  \(D_\mu\partial^\mu\omega^*=0\) and
  \(\partial_\mu D^\mu\omega=0\), in their adjointly contracted forms,
  gives \(\partial_\mu j^\mu_{\rm gh}=0\).  The generator
  \(N_{\rm gh}=\int d^3x\,j^0_{\rm gh}\) obeys the graded canonical
  relations
  \[
    [N_{\rm gh},\omega]=\omega,\qquad
    [N_{\rm gh},\omega^*]=-\omega^*,
  \]
  after the conventional factor of \(i\) is assigned to the generator.
  Thus its eigenvalue is ghost number.

  \SupplementarySolution{12}{The Free QED BRST Charge}
  \begin{enumerate}
    \item[(a)] The gauge-invariant part changes by an ordinary gauge
      transformation.  In Lorenz gauge, the variation of the gauge-fixing
      term is canceled by that of the ghosts after one integration by
      parts.  A convenient free-theory Noether current, differing from
      other choices only by improvements and equations of motion, is
      \[
        j_Q^\mu=-(\partial_\nu A^\nu)\partial^\mu\omega
          +F^{\mu\nu}\partial_\nu\omega .
      \]
      The free equations give \(\partial_\mu j_Q^\mu=0\), and
      \(Q=\int d^3x\,j_Q^0\) generates the stated graded transformations.
    \item[(b)] Inserting the standard massless mode expansions and using
      the canonical brackets gives, up to a common normalization and a
      harmless overall sign,
      \[
        Q=\int\frac{d^3p}{(2\pi)^3\,2|\mathbf p|}
        \left[
          \omega^\dagger(\mathbf p)\,p^\mu a_\mu(\mathbf p)
          +p^\mu a_\mu^\dagger(\mathbf p)\,\omega(\mathbf p)
        \right].
      \]
      Hermiticity fixes the relative coefficient of the two terms.
    \item[(c)] Taking graded commutators with this expression yields
      \[
        [Q,a_\mu]=-p_\mu\omega,\quad
        [Q,a_\mu^\dagger]=p_\mu\omega^\dagger,\quad
        \{Q,\omega^*\}=p\cdot a,\quad
        \{Q,\omega^{*\dagger}\}=p\cdot a^\dagger,
      \]
      and \(Q\) anticommutes with both \(\omega\) and
      \(\omega^\dagger\).  Hence a state
      \(\epsilon^\mu a_\mu^\dagger\ket0\) is closed when
      \(p\cdot\epsilon=0\), while \(\epsilon^\mu\sim
      \epsilon^\mu+\alpha p^\mu\) differs by an exact state.  For null
      \(p\), the quotient has the two transverse photon polarizations.
  \end{enumerate}

  \SupplementarySolution{13}{The BV Antibracket and Master Equation}
  \begin{enumerate}
    \item[(a)]
    Directly from Eq.~\eqref{eq:15.9.13},
    \[
      (\chi^n,\chi_m^\ddagger)=\delta^n_m,\qquad
      (\chi^n,\chi^m)=0,\qquad
      (\chi_n^\ddagger,\chi_m^\ddagger)=0,
    \]
    and
    \((\chi_m^\ddagger,\chi^n)=-\delta^n_m\).
    The bracket has parity one and raises ghost number by one, consistent with
    \(\epsilon(\chi_n^\ddagger)=\epsilon_n+1\) and
    \(\operatorname{gh}(\chi_n^\ddagger)=-1-\operatorname{gh}(\chi^n)\).

    The product rule gives
    \[
      (\chi^n\chi^m,\chi_k^\ddagger)
        =\chi^n\delta^m_k
         +(-1)^{\epsilon_m\epsilon_k}\delta^n_k\chi^m.
    \]
    The sign arises when the derivative associated with
    \(\chi_k^\ddagger\), whose bracket action has parity
    \(\epsilon_k\), passes through \(\chi^m\).  For two bosonic fields both
    terms have plus signs; for two fermionic fields the exchange term has a
    minus sign, as an explicit right-derivative calculation confirms.
    \item[(b)]
    Repeated use of the left- and right-derivative product rules gives
    \[
      (F,GH)=(F,G)H
        +(-1)^{(\epsilon_F+1)\epsilon_G}G(F,H)
    \]
    and
    \[
      (FG,H)=F(G,H)
        +(-1)^{\epsilon_G(\epsilon_H+1)}(F,H)G.
    \]
    Exchange \(F\) and \(G\) in the defining expression.  Moving the two
    derivatives and their results back to the original order contributes
    \((-1)^{(\epsilon_F+1)(\epsilon_G+1)}\), while exchanging the two terms
    in the odd symplectic pairing contributes an additional minus sign.
    Hence
    \[
      (F,G)=-(-1)^{(\epsilon_F+1)(\epsilon_G+1)}(G,F).
    \]
    In particular the bracket is symmetric for two even functionals and
    antisymmetric in all the other parity combinations.
    \item[(c)]
    Expanding
    \((F,(G,H))\) produces terms with a second derivative of \(G\), terms
    with a second derivative of \(H\), and, if \(E\) varies, derivatives of
    \(E\).  For constant \(E^{AB}\), graded symmetry of second derivatives
    and graded antisymmetry of \(E\) make the first two classes cancel in
    \[
      (-1)^{(\epsilon_F+1)(\epsilon_H+1)}(F,(G,H))
        +\text{cyclic permutations}.
    \]
    This proves the Jacobi identity in Darboux coordinates.

    If \(E^{AB}=E^{AB}(z)\), the surviving coefficient is the graded cyclic
    sum of
    \[
      E^{AD}\partial_D E^{BC}.
    \]
    Jacobi then requires this odd Schouten bracket \([E,E]\) to vanish.  A
    nondegenerate \(E\) obeys the condition exactly when its inverse odd
    two-form is closed.  Darboux coordinates make that form, and hence
    \(E^{AB}\), locally constant.
    \item[(d)]
    Write \(S=S_0+S_1+S_2+\cdots\) by antifield number.  The term
    \(2(S_0,S_1)=0\) is
    \[
      \frac{\delta S_0}{\delta\phi^r}R^r{}_a\omega^a=0,
    \]
    so every \(R_a\) generates a gauge symmetry.  The coefficient of
    \(\phi_r^\ddagger\omega^b\omega^c\) at the next order is
    \[
      R_b^s\partial_sR_c^r-R_c^s\partial_sR_b^r
        +C^a{}_{bc}R_a^r=0,
    \]
    up to the sign chosen in the definition of \(C\); this is closure.  The
    coefficient of
    \(\omega_a^\ddagger\omega^b\omega^c\omega^d\) is
    \[
      R_{[b}^r\partial_rC^a{}_{cd]}
        -C^e{}_{[bc}C^a{}_{d]e}=0,
    \]
    the field-dependent Jacobi identity.

    For an open algebra, closure and Jacobi hold only modulo the equations
    \(\delta S_0/\delta\phi^r=0\).  The ellipsis must then include terms
    quadratic and higher in antifields, such as
    \(\phi_r^\ddagger\phi_s^\ddagger
    M^{rs}{}_{ab}\omega^a\omega^b\).  Their antibrackets with \(S_0\)
    supply precisely the equation-of-motion terms and all higher consistency
    conditions.
    \item[(e)]
    Add the nonminimal term
    \(-\int d^4x\,h_a\omega_a^{*\ddagger}\) to \(S_{\min}\).  From the
    gauge fermion in the question,
    \[
      A_a^{\ddagger\mu}=\partial^\mu\omega_a^*,\quad
      \omega_a^{*\ddagger}
        =-\left(\partial\cdot A_a+\frac{\xi}{2}h_a\right),\quad
      h_a^\ddagger=-\frac{\xi}{2}\omega_a^*,\quad
      \omega_a^\ddagger=0.
    \]
    Substitution gives
    \[
      S_\Psi=S_{\rm YM}+\int d^4x\left[
        -(\partial_\mu\omega_a^*)(D^\mu\omega)_a
        +h_a\partial_\mu A_a^\mu+\frac{\xi}{2}h_ah_a
      \right],
    \]
    where the first sign follows by reordering the two odd factors
    \((D_\mu\omega)_a\) and \(\partial^\mu\omega_a^*\).
    The auxiliary equation \(h_a=-\partial\cdot A_a/\xi\) yields
    \[
      S_\Psi=S_{\rm YM}+\int d^4x\left[
        -(\partial_\mu\omega_a^*)(D^\mu\omega)_a
        -\frac1{2\xi}(\partial_\mu A_a^\mu)^2\right],
    \]
    the Lorenz-gauge Faddeev--Popov action.
  \end{enumerate}

  \SupplementarySolution{14}{Minimal and Quantum BV Constructions}
  \begin{enumerate}
    \item[(a)]
    With the ordering convention of Eq.~\eqref{eq:15.9.1}, the minimal action
    is
    \[
      S_{\min}=S_{\rm YM}
        +\int d^4x\left[
          (D_\mu\omega)_a A_a^{\ddagger\mu}
          -\frac12C_{abc}\omega_b\omega_c\omega_a^\ddagger
        \right].
    \]
    Here
    \[
      \operatorname{gh}(A^\ddagger)=-1,\quad
      \operatorname{gh}(\omega^\ddagger)=-2,
    \]
    while \(A^\ddagger\) is odd and \(\omega^\ddagger\) is even.  Every term
    in \(S_{\min}\) is therefore even and has ghost number zero.

    Sort \((S_{\min},S_{\min})\) by antifield number.  The term with no
    antifield is twice the gauge variation of \(S_{\rm YM}\), hence vanishes.
    Terms linear in \(A^\ddagger\) state
    \(s^2A_\mu=0\); terms linear in \(\omega^\ddagger\) state
    \(s^2\omega=0\).  The former is closure of the Yang--Mills gauge
    transformations, and the latter is the Jacobi identity.  These are
    exactly the calculations in Supplementary Exercise 15, so all
    coefficients vanish and \((S_{\min},S_{\min})=0\).
    \item[(b)]
    Because \(\Psi\) depends only on fields,
    \[
      (\chi^{\prime n},\chi^{\prime m})=0.
    \]
    Also,
    \[
      (\chi^{\prime n},\chi_m^{\ddagger\prime})
        =(\chi^n,\chi_m^\ddagger)
         -\left(\chi^n,\frac{\delta\Psi}{\delta\chi^m}\right)
        =\delta^n_m.
    \]
    Finally,
    \[
      (\chi_n^{\ddagger\prime},\chi_m^{\ddagger\prime})
        =-\frac{\delta^2\Psi}{\delta\chi^n\delta\chi^m}
         +\hbox{the graded-transposed second derivative}=0.
    \]
    The cancellation is graded symmetry of the Hessian of the fermionic
    functional.  Thus all fundamental brackets are preserved and the shift is
    anticanonical.  Setting the shifted antifields to zero immediately gives
    \[
      \chi_n^\ddagger=\frac{\delta\Psi}{\delta\chi^n},
    \]
    which is the BV gauge-fixing Lagrangian submanifold.
    \item[(c)]
    The compatibility of \(\Delta\) with the antibracket, the Jacobi identity,
    and \(\Delta^2=0\) give
    \[
      \sigma^2F
        =\frac12\left(F,(S,S)-2i\Delta S\right)
    \]
    in the right/left derivative convention of Section~15.9.  Therefore the
    quantum master equation
    \[
      (S,S)-2i\Delta S=0
    \]
    implies \(\sigma^2=0\).  Equivalently, the weighted measure
    \(e^{iS}[d\chi]\) has vanishing quantum BRST divergence.

    A quantum observable is a cohomology class of this differential:
    \[
      \sigma O=0,\qquad O\sim O+\sigma K.
    \]
    In explicit form its closure condition is
    \[
      (O,S)=i\Delta O.
    \]
    This is also the condition that makes its expectation value independent
    of the gauge-fixing Lagrangian submanifold.
  \end{enumerate}

  \SupplementarySolution{15}{Yang--Mills as a Consistent Deformation}
  Let \(f^a_{\mu\nu}=\partial_\mu A^a_\nu-\partial_\nu A^a_\mu\).
  A representative first-order deformation is
  \[
    \begin{split}
    S_1=\int d^4x\bigg[
      &-\frac12 C_{abc}f_a^{\mu\nu}A^b_\mu A^c_\nu
       +C^a{}_{bc}A_a^{\ddagger\mu}A^b_\mu\omega^c\\
      &-\frac12C^a{}_{bc}\omega_a^\ddagger
        \omega^b\omega^c\bigg].
    \end{split}
  \]
  The equation \((S_0,S_1)=0\) fixes the relative coefficients: its
  antifield-free part is the cubic vertex, while its antifield terms deform
  the Abelian gauge law and ghost transformation.  Poincaré invariance and
  integration by parts make \(C_{abc}\) totally antisymmetric with respect
  to the kinetic inner product.

  At order \(g^2\),
  \[
    (S_0,S_2)+\frac12(S_1,S_1)=0.
  \]
  The term in \((S_1,S_1)\) proportional to
  \(\omega^\ddagger\omega^3\) cannot be removed by an antifield-free
  \(S_2\); its coefficient is
  \(C^a{}_{e[b}C^e{}_{cd]}\).  Thus Jacobi is required.  When it holds, the
  remaining term is cancelled by
  \[
    S_2=-\frac14\int d^4x\,
      C^a{}_{bc}C_{ade}
      A_b^\mu A_c^\nu A^d_\mu A^e_\nu .
  \]
  Combining the quadratic, cubic, and quartic pieces gives
  \(-\tfrac14F^a_{\mu\nu}F_a^{\mu\nu}\) with
  \(F^a_{\mu\nu}=f^a_{\mu\nu}
  +gC^a{}_{bc}A^b_\mu A^c_\nu\): the unique nontrivial deformation is
  Yang--Mills theory, up to field redefinitions and strictly
  gauge-invariant additions.

\end{enumerate}
